\documentclass[../cours_tle.tex]{subfiles}
\begin{document}
\setcounter{chapter}{21}
\chapter{Espérance, variance et sommes de variables}
\label{ch:22}
\begin{objectif}
Spécialité. Démontrer les règles de calcul et repérer l'effet des dépendances. Prérequis : chapitre 21.
\end{objectif}
\section{Une variable aléatoire est une fonction}
Sur un univers fini, $X:\Omega\to\R$ associe un nombre à chaque issue. Sa loi regroupe les issues donnant la même valeur.
Son espérance est $E(X)=\sum_x xP(X=x)=\sum_{\omega}X(\omega)P(\{\omega\})$.
Sa variance est $V(X)=E((X-E(X))^2)$, son écart-type $\sigma(X)=\sqrt{V(X)}$.
L'espérance conserve l'unité de $X$, la variance possède l'unité au carré.
\section{Linéarité sans indépendance}
\begin{theoreme}
$E(aX+bY+c)=aE(X)+bE(Y)+c$, quelles que soient les dépendances.
\end{theoreme}
\begin{demonstration}
Sommer $aX(\omega)+bY(\omega)+c$ avec les poids $P(\{\omega\})$ et distribuer la somme. Les poids totalisent $1$.
\end{demonstration}
Ainsi l'espérance d'un nombre de succès est la somme des probabilités de succès, même si les succès sont dépendants.
Développer le carré autour de la moyenne donne la formule de Huygens :
$V(X)=E(X^2)-E(X)^2$. Elle implique $V(aX+b)=a^2V(X)$.
Si $V(X)=0$, $X$ vaut sa moyenne sur toutes les issues de probabilité non nulle, car une somme de carrés pondérés ne peut s'annuler autrement.
\section{Indépendance de variables et covariance}
Deux variables finies sont indépendantes si, pour toutes leurs valeurs $x,y$,
$P(X=x,Y=y)=P(X=x)P(Y=y)$.
\begin{demonstration}
Sous cette hypothèse,
$E(XY)=\sum_{x,y}xyP(X=x)P(Y=y)=(\sum_xxP(X=x))(\sum_yyP(Y=y))$.
\end{demonstration}
Définir $\operatorname{Cov}(X,Y)=E(XY)-E(X)E(Y)$ donne, par développement,
$V(X+Y)=V(X)+V(Y)+2\operatorname{Cov}(X,Y)$.
Pour des variables indépendantes, la covariance est nulle, donc les variances s'ajoutent. Une covariance nulle n'implique pas l'indépendance : si $X$ est uniforme sur $\{-1,0,1\}$ et $Y=X^2$, $E(XY)=E(X)=0$, mais $Y$ dépend de $X$.
\section{Bernoulli, binomiale et moyenne d'échantillon}
Pour $B\sim\mathcal B(1,p)$, $B^2=B$, donc $E(B)=p$, $V(B)=p-p^2=p(1-p)$.
Si $X=B_1+\cdots+B_n$ avec Bernoulli indépendantes de même paramètre,
$E(X)=np$ et $V(X)=np(1-p)$.

Un échantillon de taille $n$ d'une loi est ici une famille de variables indépendantes de cette même loi. Si elles ont moyenne $\mu$ et variance $\sigma^2$, leur moyenne $M_n$ vérifie
\[
E(M_n)=\mu,\qquad V(M_n)=\frac{\sigma^2}{n},\qquad \sigma(M_n)=\frac{\sigma}{\sqrt n}.
\]
La première formule vient de la linéarité ; la seconde de l'indépendance et du facteur $1/n^2$ appliqué à la variance de la somme.
\section{Exemples et erreurs à détecter}
La somme de deux dés indépendants a espérance $7$, même sans écrire toute sa loi.
Pour un dé uniforme, $E(D^2)=91/6$, donc $V(D)=35/12$ et la variance de la somme vaut $35/6$.
Si l'on utilise deux fois le même résultat, la variable $2D$ a aussi espérance $7$ mais variance $4V(D)=35/3$ : la dépendance change la dispersion.
Pour une moyenne, quadrupler $n$ divise l'écart-type par $2$, pas par $4$.
\begin{finchapitre}
\section{Utilité et suite du parcours}
Les variances et covariances forment une géométrie des variables aléatoires. Elles conduisent aux projections statistiques, aux moindres carrés et à l'étude des estimateurs.
\begin{retenir}
L'espérance est toujours linéaire. L'addition des variances exige une justification de l'absence de termes de covariance.
\end{retenir}
\end{finchapitre}
\end{document}
